%! Fundamental Theorem of Algebra and Complex Zeros — Unit 3, Lesson 3.5 — Warm-Up (Answer Key)
\documentclass[10pt]{article}
\usepackage{algebra2-article}
\usepackage{algebra2-key}   % pulls in -boxes; do NOT also load -boxes
\newcommand{\TallMath}[1]{$\displaystyle #1\rule[-1.4em]{0pt}{3.2em}$}

\begin{document}

\pageheader{Unit 3, Lesson 3.5}{Warm-Up --- Answer Key}
\namedateperiod

\vspace{0.10in}

\begin{notesbox}{Getting Ready: Where the Missing Solutions Went}
\small Yesterday a cubic gave you \emph{one} solution and then stalled. Three things you already know
--- one from Lesson 2.4, one from Lesson 2.5, and one from yesterday --- are enough to finish it.

\vspace{5pt}
\textbf{1. Conjugates multiply to something real} (Lesson 2.4). Recall $i^2 = $ \ans{$-1$}

\vspace{3pt}
Multiply and simplify: \quad $(x - 3i)(x + 3i) = $ \ans{$x^2 - 9i^2 = x^2 + 9$}

\vspace{3pt}
Is your answer a \emph{real} polynomial --- no $i$ left anywhere? \ans{yes} \quad
Where did the $i$'s go? \ans{the $\pm 3ix$ middle terms cancelled}

\vspace{6pt}
\textbf{2. Two ways to finish a quadratic} (Lessons 2.3 and 2.5).

\vspace{3pt}
(a) Square roots: \quad $x^2 + 4 = 0 \;\Rightarrow\; x^2 = $ \ans{$-4$} $\;\Rightarrow\; x = $
\ans{$\pm 2i$}

\vspace{3pt}
(b) Quadratic formula: \quad $x^2 - 2x + 5 = 0$. \quad Discriminant $b^2-4ac = $ \ans{$4 - 20 = -16$}

\vspace{2pt}
\hspace{1.2em} So there are \ans{no (zero)} real solutions, and $x = $ \ans{$\frac{2 \pm 4i}{2} = 1 \pm 2i$}

\vspace{3pt}
(c) Look at your two answers in (a), and your two answers in (b). What do the two solutions in each
pair have in common? \ans{same real part, opposite imaginary parts --- they are conjugates}

\vspace{6pt}
\textbf{3. Unfinished business.} Yesterday your group solved $x^3 - x^2 + 4x - 4 = 0$. The candidate
$r = 1$ worked, and dividing gave the depressed polynomial $x^2 + 4$ --- where you stopped.

\vspace{3pt}
Completely factored: $x^3 - x^2 + 4x - 4 = $ \ans{$(x-1)(x^2+4)$}

\vspace{3pt}
Set each factor equal to zero and solve \emph{both} of them, using item 2(a) for the second:

\vspace{3pt}
\hspace{1.2em} From the first factor: $x = $ \ans{$1$} \qquad
From the second: $x = $ \ans{$\pm 2i$}

\vspace{3pt}
All the solutions: \ans{$x = 1,\ 2i,\ -2i$}

\vspace{6pt}
\textbf{4. The count.} Yesterday this cubic seemed to have $1$ solution.

\vspace{2pt}
How many does it have now? \ans{$3$} \quad What is the degree of the polynomial? \ans{$3$}

\vspace{2pt}
Write down what you think that coincidence means: \ans{the degree may tell you how many solutions
there are --- if imaginary ones count}

\end{notesbox}

\begin{teachernote}
This page is the lesson in miniature, and nothing on it is new. Item~1 rebuilds the conjugate product
from Lesson~2.4 --- the mechanism that lets a polynomial with \emph{real} coefficients own an
imaginary zero. Item~2 rebuilds the two finishing moves for a quadratic and, in part (c), gets
students to \emph{notice} the pairing before it is named. Item~3 is Lesson~3.4's Tier E cliffhanger,
finished in two lines now that $\pm2i$ is allowed as an answer.

\vspace{3pt}
\textbf{Debrief items 1 and 3 together.} Ask: \textbf{``In item 3 you got $2i$ and $-2i$. In item 2(b)
you got $1+2i$ and $1-2i$. Is that a coincidence?''} Take answers --- do not name the Complex
Conjugate Root Theorem yet. Then point back at item~1: the reason the pair has to travel together is
that only the pair multiplies back to something real. That is the first minute of the notes, and it
lands much harder if students have seen the pattern three times on this page first.

\vspace{3pt}
Item~4 is the hinge of the whole lesson: degree $3$, three solutions. Resist confirming it. Students
who write ``the degree tells you how many answers there are'' have just stated the corollary of the
Fundamental Theorem of Algebra on their own, and the notes should read as confirmation of their
conjecture, not as an announcement. Common wrong turns to watch for: $x = 2i$ only in 2(a) (lost the
$\pm$), $x^2 - 9$ in item~1 (forgot $i^2 = -1$), and ``no solution'' still written in 2(b) --- that
last one is exactly the habit the period replaces.
\end{teachernote}

\end{document}
